% ==================================================================
% id:       cond-024
% topic:    הסתברות מותנית ואי-תלות
% subtopic: מאורעות בלתי תלויים
% solution: full
% ==================================================================
יהיו $A$ ו־$B$ מאורעות בלתי תלויים במרחב הסתברות, כך ש־$\PP(A)=0.5$ ו־$\PP(B)=0.6$.
\begin{enumerate}
\item מה ההסתברות שיקרה בדיוק אחד מהם?
\hint{״בדיוק אחד״ הוא איחוד זר של שני מאורעות. בכל אחד מהם מופיע משלים.}
\begin{solbox}
נסמן $E$ = ״קרה בדיוק אחד מהמאורעות $A,B$״. מחפשים את $\PP(E)$.

״בדיוק אחד״ פירושו: $A$ קרה ו־$B$ לא, או $B$ קרה ו־$A$ לא. כלומר $E=(A\cap B^C)\cup(A^C\cap B)$, ושני המאורעות האלה זרים. לכן
\[
\PP(E)=\PP(A\cap B^C)+\PP(A^C\cap B).
\]
אם $A$ ו־$B$ בלתי תלויים, גם $A$ ו־$B^C$ בלתי תלויים, וכך גם $A^C$ ו־$B$. לכן
\[
\PP(E)=\PP(A)\,\PP(B^C)+\PP(A^C)\,\PP(B)=0.5\cdot0.4+0.5\cdot0.6=0.2+0.3=\ans{0.5}.
\]
\end{solbox}
\item מה ההסתברות שיקרה לפחות אחד מהם?
\begin{solbox}
$\PP(A\cup B)=\PP(A)+\PP(B)-\PP(A)\PP(B)=0.5+0.6-0.3=\ans{0.8}$.
\end{solbox}
\item מה ההסתברות שלא יקרה אף אחד מהם?
\begin{solbox}
זה המשלים של סעיף (ב): $1-0.8=\ans{0.2}$. אפשר גם ישירות: $\PP(A^C\cap B^C)=0.5\cdot0.4=0.2$.
\end{solbox}
\end{enumerate}
